% The final report should read like a 6-8 page NLP/ML conference paper and use the included report_template.tex.
\documentclass[10pt]{article}
\usepackage[margin=0.72in]{geometry}
\usepackage[T1]{fontenc}
\usepackage{lmodern}
\usepackage{microtype}
\usepackage{booktabs}
\usepackage{graphicx}
\usepackage{amsmath,amssymb}
\usepackage{xcolor}
\usepackage{hyperref}
\usepackage{enumitem}
\definecolor{hkustblue}{HTML}{003B71}
\hypersetup{colorlinks=true,urlcolor=hkustblue,citecolor=hkustblue}
\setlength{\parindent}{0pt}
\setlength{\parskip}{4pt}

\title{\vspace{-1.5em}\textbf{Your Project Title}}
\author{Student Name (Student ID) \and Student Name (Student ID)}
\date{COMP 4221 Natural Language Processing --- Fall 2026}

\begin{document}
\maketitle

\begin{abstract}
In no more than 250 words, state the problem, approach, data, principal result, and main conclusion. Replace all instructional text before submission.
\end{abstract}

\section{Introduction}
Define the task and motivation. State the research question or claim, summarize the approach, and list the main findings or contributions.

\section{Related Work}
Position the project against the most relevant literature. Cite sources, including reused models, datasets, code, and evaluation tools.

\section{Data and Task Formulation}
Describe data provenance, licensing, preprocessing, splits, task definition, and important limitations. Explain any privacy or ethical considerations.

\section{Approach}
Describe the proposed method or analysis precisely. Include equations, an algorithm, or a system figure when useful. Explain design choices and identify what is reused versus implemented by the team.

\section{Experimental Setup}
Specify baselines, metrics, hyperparameters, model versions, random seeds, and compute. Explain how the evaluation answers the research question.

\section{Results}
Report the main quantitative comparisons in tables or figures. State uncertainty or variation when appropriate. Do not report numbers without interpreting them.

\section{Analysis}
Provide qualitative examples, error categories, ablations, or subgroup analysis. Explain when and why the approach succeeds or fails.

\section{Limitations and Responsible Use}
Discuss limitations, possible harms, dataset or model biases, and situations in which the conclusions may not generalize.

\section{Conclusion}
Answer the research question using the evidence presented. Suggest a small number of concrete next steps.

\section*{Team Contributions and Use of AI Tools}
List each member's contributions and approximate percentage. Disclose material use of generative AI or other automated tools for coding, analysis, or writing, and describe how outputs were checked.

\section*{Reproducibility Statement}
State where the code is included, the main command to reproduce results, expected runtime, hardware or API requirements, and any unavailable resources.

\bibliographystyle{plain}
\bibliography{references}

\appendix
\section{Additional Results}
Optional material may appear here. The appendix and references do not count toward the 6--8 page main-text limit.

\end{document}
